\documentclass[11pt]{article}
\usepackage[margin=0.9in]{geometry}
\usepackage{amsmath,amssymb,amsthm,mathtools,hyperref}
\hypersetup{hidelinks,pdftitle={No fixed separable polynomial multiplier for canonical Bell slacks}}
\newtheorem{theorem}{Theorem}
\newtheorem{corollary}[theorem]{Corollary}
\newcommand{\R}{\mathbb R}
\title{No fixed separable polynomial multiplier\\for canonical Bell slacks}
\author{}\date{}
\begin{document}
\maketitle\vspace{-2em}

A singular limit of doubly tilted CHSH produces the Motzkin
polynomial~\cite{Chaturvedi}. We use an explicit strictly valid source
family and a polynomial four-component state to exclude every fixed
separable polynomial multiplier. The threshold obtained here is
existential; the multiplier may be chosen arbitrarily before taking
the limit.

Put $D_u=1+u^2$, $D_v=1+v^2$, and let
$\sigma_x=\left(\begin{smallmatrix}0&1\\1&0\end{smallmatrix}\right)$,
$\sigma_z=\operatorname{diag}(1,-1)$. On two qubits choose
\[
 A_0=\sigma_z\otimes I,\quad
 A_1=\frac{(1-u^2)\sigma_z+2u\sigma_x}{D_u}\otimes I,
 \qquad
 B_0=I\otimes\sigma_z,\quad
 B_1=I\otimes\frac{(1-v^2)\sigma_z+2v\sigma_x}{D_v}.
\]
For $0<T\le1/100$, set $\eta=1-T$, $w=4-2T+4T^3$, and define
\begin{equation}
 S=A_0B_0-A_0B_1+A_1B_0+A_1B_1,\qquad
 F_T(u,v;\psi)=D_uD_v\,\psi^T
       [wI-S-\eta(A_1+B_0)]\psi .                 \label{eq:parent}
\end{equation}
Here $\psi=(W,X,Y,Z)^T$ is an arbitrary unnormalised real state.
Thus $F_T\in\R[u,v,W,X,Y,Z]$ is defined without denominators.
In the canonical source convention of~\cite{Canonical}, its coefficients
$c(T)=(d,\lambda,n,m,C,A,D,\alpha,\beta)$ are exactly
\begin{equation}
 d=2T-4T^3,\qquad \lambda=m=A=T^3,\qquad
 n=D=\alpha=1-T+T^3,\qquad C=\beta=(1-T)/2+T^3.
                                                        \label{eq:source}
\end{equation}
This follows by expanding
$d=4-w$, $\lambda=m=A=(w-2-2\eta)/4$,
$n=D=\alpha=(w-2+2\eta)/4$, and $C=\beta=(w-2)/4$.

\begin{theorem}\label{thm:main}
Every source \eqref{eq:source} is strictly valid for all states and
all projective measurement angles, and all nine source coefficients
are positive. For every fixed pair of nonzero polynomials
$p\in\R[u]$, $q\in\R[v]$ nonnegative on the whole real line, there is
$\delta_{p,q}>0$ such that
\begin{equation}
 p(u)q(v)F_T(u,v;W,X,Y,Z)
 \quad\text{is not a polynomial sum of squares}
 \qquad(0<T<\min\{\delta_{p,q},1/100\}).            \label{eq:main}
\end{equation}
In particular, the rational sources $T=1/n^2$ eventually exclude every
prescribed pair $p,q$.
\end{theorem}

\begin{proof}
\emph{Validity.}
Swap Alice's inputs and rotate her first observable to $\sigma_z$.
For real $c,r,e,s$ with $c^2+r^2=e^2+s^2=1$, the Bell operator becomes
\begin{align*}
 H_T={}&(1+c+e-ce)\,\sigma_z\otimes\sigma_z
 +(1-c)s\,\sigma_z\otimes\sigma_x
 +(1-e)r\,\sigma_x\otimes\sigma_z\\
 &-rs\,\sigma_x\otimes\sigma_x
 +(1-T)(\sigma_z\otimes I+I\otimes\sigma_z).
\end{align*}
Write $a=1-c$, $b=1-e$, so $a,b\in[0,2]$. Direct expansion gives
\begin{align}
 \det[(4-2T+\rho)I-H_T]
 &=\rho P_T(\rho;a,b)+8abR_T(a,b),                 \label{eq:det}\\
 P_T&=\rho^3+(16-8T)\rho^2+(84-88T+20T^2)\rho\notag\\
 &\quad+144-240T+112T^2-16T^3+8ab(1-T)^2,\notag\\
 R_T&=(3+2T-T^2)(a+b)-2ab-18T+12T^2-2T^3.\notag
\end{align}
At $\rho=4T^3$, $P_T\ge144-240T$. If $R_T\ge0$, the determinant
is positive. Otherwise let $z=a+b$. Since $ab\le a+b$,
$R_T\ge z-18T$ and hence $z<18T<2$. Now $2ab\le z^2/2\le z$
gives $R_T\ge2z-18T$, so
\[
 8abR_T\ge2z^2R_T\ge4z^3-36Tz^2\ge-432T^3;
 \qquad
 4z^3-36Tz^2+432T^3=4(z-6T)^2(z+3T).
\]
Consequently
$\det(wI-H_T)\ge(144-960T)T^3>0$ in this case also.
At $c=e=1$, the four eigenvalues of $wI-H_T$ are
\[
 4T^3,\quad6-2T+4T^3,\quad6-2T+4T^3,\quad4-4T+4T^3.
\]
They are positive. The angle torus is connected and the determinant
never vanishes there, so its inertia is constant. This proves strict
positivity for both independent sine signs, including commuting
angles and projective endpoints. Real and imaginary state parts, and
the standard two-observable qubit reduction, give validity on
$Q_{222}$. Positivity of every coefficient in \eqref{eq:source} is literal.

\emph{A polynomial full-state test.}
Put $T=t^2$, $u=-tx$, $v=ty$ and take
\begin{equation}
 \psi_t(x,y)=
 \begin{pmatrix}
 1+\frac23t^4x^2y^2\\
 -\frac13t^3x^2y\\
 -tx+\frac23t^3xy^2\\
 -t^2xy-\frac23t^4x^3y
 \end{pmatrix}.                                      \label{eq:state}
\end{equation}
The choice cancels the lower-order state contributions as follows.
In the standard frame use half-angles $tx,ty$ and
$\xi=(1,\mu t^3,\nu t^3,h t^2)^T$. The coefficient of $t^4$ in
$\xi^T(wI-H_{t^2})\xi$ is $4(h+xy)^2$. At $h=-xy$, its coefficient
of $t^6$ is
\[
 6\left(\mu-\frac23x^2y\right)^2
 +6\left(\nu-\frac23xy^2\right)^2
 +\frac43\{M(x,y)+2\},\qquad
 M=x^4y^2+x^2y^4-3x^2y^2+1.
\]
Choose the two displayed state amplitudes to cancel the first two squares.
Multiplication by
$\left(\begin{smallmatrix}1&-u\\u&1\end{smallmatrix}\right)\otimes I$
then gives exactly \eqref{eq:state}. This polynomial rotation numerator
and angle clearing multiply the expectation by $D_u^2D_v$, with constant
term one. Direct polynomial expansion therefore gives
\begin{equation}
 F_{t^2}(-tx,ty;\psi_t)
 =\frac43t^6\{M(x,y)+2\}+t^8R(t,x,y),
 \qquad R\in\mathbb Q[t,x,y].                        \label{eq:pullback}
\end{equation}
Its total degree in $(x,y)$ is at most twelve.

\emph{The fixed multiplier.}
Write the first nonzero Taylor terms as
$p(u)=a_0u^{2j}+O(u^{2j+1})$ and
$q(v)=b_0v^{2k}+O(v^{2k+1})$. Global nonnegativity gives
$a_0,b_0>0$. If the conclusion failed, there would be a sequence
$t_i\downarrow0$ for which $pqF_{t_i^2}$ is SOS. Substitute
\eqref{eq:state} and divide by $t_i^{6+2j+2k}>0$.
The resulting SOS polynomials in the independent variables $x,y$
have total degree at most $12+\deg p+\deg q$ and converge coefficientwise to
\begin{equation}
 \frac{4a_0b_0}{3}\,x^{2j}y^{2k}\{M(x,y)+2\}.          \label{eq:limit}
\end{equation}
The SOS cone of fixed degree is closed. Indeed, let $z(x,y)$ contain
all monomials through the corresponding root degree. The matrix
$B=\int_{[-1,1]^2}zz^T\,dx\,dy$ is positive definite.
If $Q_i\succeq0$ represents the $i$th polynomial, coefficient
convergence bounds its integral $\operatorname{Tr}(BQ_i)$ and hence
$\operatorname{Tr}Q_i$. A convergent subsequence of Gram matrices
then represents the limit by a positive semidefinite matrix.

Thus \eqref{eq:limit} would be SOS. A square monomial can be cancelled
from an SOS: evaluating a sum of squares at $x=0$ forces each root
to be divisible by $x$, and repeated division removes $x^{2j}$;
the same argument removes $y^{2k}$. After removing the positive scalar,
$M+2$ would be SOS. Every root would have degree at most three.
Since $M+2=3$ on either coordinate axis, each root is constant on
both axes, hence has the form
$a+bxy+cx^2y+dxy^2$. The coefficient of $x^2y^2$ in a sum of these
squares is nonnegative, whereas it is $-3$ in $M+2$.
This contradiction rules out every sequence approaching $T=0$ and
proves the required sufficiently-small-$T$ statement.
\end{proof}

\begin{corollary}\label{cor:completion}
No fixed nonzero polynomial $q\in\R[v]$ nonnegative on $\R$ permits
the canonical two-LMI polynomial completion for every valid source.
More precisely, for each such $q$, all sufficiently small sources
\eqref{eq:source} admit no polynomial completion matrix.
\end{corollary}
\begin{proof}
Use the canonical completion matrices $G_{q,Z}(v),H_{q,Z}(v)$
of~\cite[Sections 1 and 4]{Canonical}: the prescribed source blocks
are multiplied by $q$, and the off-diagonal completion is the
polynomial matrix $Z$. Their exact splitting identity is
\[
 q\Phi_c(U,v;a,b)=
 \binom a{Ub}^{\!T}G_{q,Z}(v)\binom a{Ub}
 +U\binom ab^{\!T}H_{q,Z}(v)\binom ab.
\]
Here $\Phi_c$ is the canonical quadratic form in the transformed
state coordinates $(a,b)$, and $\mathcal P_c$ below is its state matrix.
If both univariate matrix polynomials were positive semidefinite
on all of $\R$, polynomial matrix-SOS factorization~\cite[Theorem 1.2]{Zalar}
would make $q\Phi_c(u^2,v;a,b)$ SOS. The exact polynomial inverse-numerator identity
from~\cite[Section 1]{Canonical} is
\[
 (L(u)\psi)^T\mathcal P_c(u^2,v)(L(u)\psi)
 =u^2(1+u^2)F_c(u,v;\psi),\qquad
 L(u)=\begin{pmatrix}
 u&0&u^2&0\\0&u&0&u^2\\u&0&-1&0\\0&u&0&-1
 \end{pmatrix}.
\]
Substitution would therefore make $q(v)u^2(1+u^2)F_T$ SOS, contrary
to Theorem~\ref{thm:main} with $p(u)=u^2(1+u^2)$.
\end{proof}

The obstruction concerns fixed separable polynomial multipliers;
it does not address source-dependent or nonseparable multipliers.
A square-root testing function already removes this particular limiting
obstruction on a half-line: the exact identity
\[
 M(x,s^2)=(x^2s^2-1)^2
 +\bigl(xs^2(s^2-1)\bigr)^2
 +2\bigl(xs(s^2-1)\bigr)^2
\]
makes $M(x,s^2)+2$ SOS. This identity supplies no universal half-line
completion or finite global semidefinite lift.

\begin{thebibliography}{3}\small
\bibitem{Chaturvedi} A.~Chaturvedi,
\emph{No Finite NPA Level Characterizes the Complete Quantum Set
in the Simplest Bell Scenario},
\href{https://arxiv.org/abs/2607.14569}{arXiv:2607.14569v1} (2026),
Sections S3 and S6--S7.
\bibitem{Canonical} A.~Chaturvedi,
\emph{Canonical Bell validity through two matrix inequalities},
\href{https://github.com/anubhav70747/q222-finite-sdp/tree/main/extracted-lifts/canonical-completion}{q222-finite-sdp, canonical-completion note}.
\bibitem{Zalar} A.~Zalar,
\emph{Matrix Fej\'er--Riesz theorem with gaps},
\href{https://arxiv.org/abs/1503.06034}{arXiv:1503.06034}, Theorem~1.2.
\end{thebibliography}
\end{document}
